\chapter{The geometric longitude selected by the finite source}
\label{gc:chapter}

The meridian profile contains more information than its first cusp
correction. For every covered fixed strict Conway tree with
independently sufficiently long signed twists, its full classical
characteristic is the preferred-longitude curve of the original
knot complement. Product junctions may have arbitrary valence and
may be adjacent. Ordered trees carrying a crossing circle at every
junction also retain their separate quantitative filling criteria.

We first establish the identity for the circle trees. Inserting a
circle around a zero-crossing box then supplies the geometric
extension to a bare product junction. Eliminating its angular
variable recovers the original stationary action with its constant,
while the ordered boundary paths recover the original knot group.
The consequence for a quantum recurrence uses the separate all-order
source theorem; the geometric statement alone supplies no new
asymptotic estimate for short fillings.

\section{The marked equations and their geometric branch}
\label{gc:equations}

Keep the literal construction of
Section~\ref{sec:circle-the-knot-construction}: one original crossing
circle at every vertex of a fixed finite ordered tree, at least one
edge, and an actual quarter-turn on every edge. The main strands must
form one knot, with no additional component closed inside a tangle.
All disk framings, endpoint orders and signed powers
$n_v=\epsilon_v-2p_v$ are retained.

Write $\mathcal E$ for the edge set. For an edge $e=vw$, let
$E_e=e^{2\pi i x_e}$ be its complex wire parameter, and put
$Q_v=e^{i\tau_v}$ at a crossing circle. These variables describe
the original cells, rather than a new diagram obtained by changing
the products. Let $M$ be the main meridian parameter. The selected
source flag has meridian eigenvalue $M^{-1}$; throughout this
chapter $L$ is the inverse eigenvalue of the preferred longitude
on that same flag. Equivalently, $(M,L)$ is the eigenvalue pair
on the other line of the determinant-one peripheral representation.

The finite source gives the rational equations
\begin{equation}\label{gc:edge}
\begin{aligned}
 &(E_e+Q_v)(E_e+Q_v^{-1})(E_e+Q_w)(E_e+Q_w^{-1})\\
 &\hspace{12mm}=(E_e-1)^2(1-E_eM^2)(1-E_e/M^2),
\end{aligned}
\end{equation}
\begin{equation}\label{gc:vertex}
 \prod_{e\ni v}\rho_{v,e}=Q_v^{n_v},
 \qquad
 \rho_{v,e}=\frac{Q_v+E_e}{1+Q_vE_e}.
\end{equation}
We use the chart on which these denominators are nonzero and the
selected critical point is nondegenerate, meaning that its Hessian
is invertible. Exponentiation alone
does not specify that point. To retain its logarithmic branch, put
\[
 b(x,\tau)=\int_0^x
       \log(2\cos(2\pi r)+2\cos\tau)\,dr-2J(x),
 \qquad J'(x)=\log(2\sin\pi x),
\]
with the original continued logarithms. The angular equations are
\begin{equation}\label{gc:first-charge}
 \sum_{e\ni v}b_\tau(x_e,\tau_v)
            +\frac{i n_v\tau_v}{2\pi}-i=0.
\end{equation}
The term $-i$ is the first Fourier charge of the actual crossing
trace. The positive charge path of the preceding geometric
completion starts at the complete augmented structure and ends
at this filled point. It identifies the source critical point
with the actual complete knot representation. The implicit-function
theorem then selects its continuation as $M$ varies near one.

\begin{theorem}[The preferred eigenvalue on the regular branch]
\label{gc:longitude-theorem}
Fix the construction above and a regular geometric branch through
an actual complete hyperbolic filling, with the logarithms selected
by \eqref{gc:first-charge}. On its continued germ,
\begin{equation}\label{gc:longitude}
 \boxed{\displaystyle
 L_{\mathrm{geom}}(M)
   =-M^{a(D)}
          \prod_{e\in\mathcal E}\frac{M^2-E_e(M)}
                                      {1-M^2E_e(M)},\qquad
 a(D)=-w(D)-\sum_vn_v .}
\end{equation}
Here $w(D)$ is the oriented writhe of the original filled diagram.
The longitude is its preferred zero-linking longitude, and the
representation is that of the original knot complement.
\end{theorem}

This theorem concerns a regular geometric branch. The long-filling
and quantitative hypotheses in the preceding chapters provide such
branches, with their stated knot and parity qualifications. They
remain the hypotheses for the corresponding volume limits; the
identification \eqref{gc:longitude} does not enlarge those analytic
claims.

\section{Why the equations describe the original exterior}
\label{gc:actual-exterior}

The attachment in Figure~\ref{fig:circle-octahedral-attachment}
is the topological starting point. Each original tree edge supplies
two octahedra. Their disk faces are sewn in the original cyclic
order, retaining the distinguished odd seam at each circle.
For a nonparabolic main meridian, both projection copies and their
white-face maps must deform. Keeping the old parabolic white maps
would give the wrong representation.

The resulting marked edge block satisfies \eqref{gc:edge}.
Its primitive auxiliary longitude has similarity $Q_v^2$,
and its meridian contribution is $\rho_{v,e}$.
The pure page identifications match both puncture generators.
At a positive half twist the actual ordered pair changes by
$(a,b)\mapsto(a^{-1}ba,a)$, preserving $ba$.
The invertible matrix $(ba)^{-1}-I$ realizes that conjugacy
near the geometric chart. The accompanying puncture flags must
also be transported: the two native ideal flags can select
opposite eigenvalues. The full generator matching, including
this exchange, is what permits gluing the actual collar.

Consequently the original primitive logarithms are
\begin{equation}\label{gc:primitive}
 U_v=\sum_{e\ni v}\log\rho_{v,e}-\epsilon_v\log Q_v,
 \qquad V_v=2\log Q_v,\qquad U_v+p_vV_v=-2\pi i.
\end{equation}
The last equality follows from the marked equation
\eqref{gc:first-charge}, since
$\exp(-2\pi b_\tau)=\rho_{v,e}$.
It retains the primitive period which disappears from
\eqref{gc:vertex}. The same ordered face and collar relations therefore construct
representations of the original augmented exterior. At the selected
filled point, $\operatorname{Im}\tau_v<0$, so $Q_v^2\ne1$.
The auxiliary longitude is noncentral and semisimple. Its commuting
filling word has eigenvalue ratio one by \eqref{gc:vertex}, and is
therefore projectively the identity. This proves descent to the
prescribed filled knot group, also on a sufficiently small nearby
germ. At $Q_v=\pm1$, an exponentiated slope equation alone would
not exclude a nontrivial parabolic filling word.

\section{From the peripheral differential to an exact identity}
\label{gc:rigidity}

There is a useful algebraic way to determine the entire longitude
without inferring it from a tangent. Put
$u=\log M$, $q_v=\log Q_v$, and
$f_e=(M^2-E_e)/(1-M^2E_e)$.
Differentiating the single edge equation gives
\begin{equation}\label{gc:edge-differential}
 du\wedge d\log f_e
 +\frac12 d\log\rho_{v,e}\wedge dq_v
 +\frac12 d\log\rho_{w,e}\wedge dq_w=0.
\end{equation}
The actual representation has the same boundary pairing:
\[
 du\wedge d\log L_{\mathrm{aug}}
       +\frac12\sum_v dU_v\wedge dq_v=0.
\]
Here $L_{\mathrm{aug}}$ uses the preferred augmented longitude.
This identity follows by pairing two infinitesimal representation
cocycles with the invariant trace form and evaluating their cup
product on the boundary of the exterior. The cochain Stokes
identity makes the total boundary evaluation zero. The factor
$1/2$ records that $U_v$ is a logarithm of an eigenvalue ratio.
Subtracting \eqref{gc:edge-differential} over all edges removes
the auxiliary terms in \eqref{gc:primitive}.

It follows that the rational squared discrepancy
\begin{equation}\label{gc:discrepancy}
 R=\frac{L_{\mathrm{aug}}^2}{\prod_e f_e^2}
\end{equation}
has no infinitesimal dependence on the auxiliary variables.
Two further facts turn that differential statement into an exact
formula. First, the generic source field has no extra algebraic
constants over $\mathbb C(M)$. Indeed its edge squareclasses are
\[
 (Q_v-1)^2(Q_w-1)^2
                 -4(M+M^{-1})^2Q_vQ_w.
\]
Each is an irreducible prime involving its own pair of vertex
variables. Their prime valuations make the squareclasses
independent, even after algebraically closing $\mathbb C(M)$.
Thus $R\in\mathbb C(M)$.

Second, at every fixed $M_0\ne0$ the group matrices and their
inverses admit a regular reconstruction for generic auxiliary
variables. This includes $M_0=\pm i$, where a particular ideal
subdivision degenerates but the character matrices remain regular.
The peripheral eigenvalue ratio and every $f_e$ are nonzero
there. Hence $R$ has neither a zero nor a pole on
$\mathbb C^*$, and therefore
\begin{equation}\label{gc:monomial}
                         R=cM^k
\end{equation}
for a nonzero constant $c$ and an integer $k$.
This regularity step is essential: a differential identity and a
single cusp derivative would not exclude a general rational
factor of $M$.

At the complete main cusp, $R(1)=1$. The already marked cusp
formula fixes $k=2a_{\mathrm{aug}}$, where
$a_{\mathrm{aug}}=-\sum_v\epsilon_v-w_\epsilon$.
The original surgery homology changes the preferred longitude by
$\mu_K^{\sum_vp_v\ell_v^2}$, with
$\ell_v\in\{0,\pm2\}$ the original circle linking numbers.
Its exponent is
\[
 a_{\mathrm{aug}}+\sum_vp_v\ell_v^2
             =-w(D)-\sum_vn_v=a(D).
\]
This proves the square of \eqref{gc:longitude}.

The remaining sign has a topological determination. A spin cover
is the double cover of oriented frames in which a full $2\pi$
rotation changes the lift's sign. On a Seifert
surface with one boundary component, a boundary-adapted tangent
frame winds $1-2g$ times relative to a global tangent frame.
The spin cover of that global frame has trivial boundary
monodromy, because the boundary is nullhomologous on the surface.
Thus the adapted frame lifts antiperiodically. At the straight
preferred cusp longitude it is the translation-invariant
horospherical frame; its positive unipotent lift therefore differs
from the actual longitude lift by a minus sign.
So $L_{\mathrm{geom}}(1)=-1$ at the complete filled point.
The two-valued holomorphic ambiguity in the squared identity is
constant, proving \eqref{gc:longitude}.
Appendix~\ref{app:geometric-characteristic} gives the marked
matrices, regularity argument and frame proof in full.

\section{Product junctions without crossing circles}
\label{gc:strict-extension}

A bare junction is a maximal consecutive product that contains no
crossing-circle operator. A binary association used to calculate its
kernel introduces hidden labels, but adds no physical quarter-turns.
The strict construction retains the original ordered products,
turns and cylindrical closure. Its real angle domain contains a
point at which every original leaf angle lies in $(0,1/2)$ and
every retained binary gate obeys
\[
                 \max(u,v)<w<u+v<1.
\]
The original endpoint test must give one knot, including any circles
closed inside a tangle. The two-leaf pure-leaf qualification and all
original parity conditions remain in force.

Let $\mathcal E_T$ denote the physical turn edges of the reduced
junction tree. Their covariant index $x_e$ is shared by the two ends
after lowering the trace pairing. When an incidence is exposed as
the ordinary output of its rooted branch, that output angle is
$1-x_e$. Put $X_e=e^{2\pi ix_e}$; this is the variable denoted $E_e$
in the preceding circle-tree formulas.

\begin{theorem}[The longitude for a strict Conway tree]
\label{gc:strict-longitude-theorem}
Fix such a strict ordered knot presentation and independently
sufficiently long signed powers, with the diagram and all its powers
fixed before the level varies. Its original complete hyperbolic
filled knot has a regular geometric main-meridian germ on which
\begin{equation}\label{gc:strict-longitude}
 L_{\mathrm{geom}}(M)
   =-M^{a(D)}\prod_{e\in\mathcal E_T}
           \frac{M^2-X_e(M)}{1-M^2X_e(M)},
 \qquad a(D)=-w(D)-\sum_i n_i.
\end{equation}
The $x_e(M)$ are selected by the original strict-tree action
continued with the physical-edge insertion
\begin{equation}\label{gc:strict-phase}
\begin{aligned}
 \Phi_t&=\Phi+\sum_{e\in\mathcal E_T}
       \{2J(x_e)-J(x_e+t)-J(x_e-t)\}
                 +\frac{\pi i a(D)}2t^2,\\
 &\hspace{32mm}M=e^{\pi it}.
\end{aligned}
\end{equation}
The preferred peripheral curves and the filled complement are those
of the original knot.
\end{theorem}

The proof begins with the framed collar in
Figure~\ref{fig:strict-zero-circle-junction}. Insert the literal
identity $R^0=I$ immediately before the ordered child product at
each bare junction, and encircle its two straight strands. Filling
that new circle along its primitive meridian restores the same
framed tangle. This operation creates no turn edge, crossing power
or preferred-longitude shear. It therefore leaves both
$\mathcal E_T$ and $a(D)$ unchanged.

There is a unique regular angular variable for each new circle.
If $x_1,\ldots,x_m$ are its incident covariant labels at the complete
strict optimizer, then $0<x_j\le1/2$, at most one label equals
$1/2$, and $\sum_jx_j>1$. With $k=\tanh(y/2)$, its zero-power
first-charge equation is
\begin{equation}\label{gc:zero-angular}
 \sum_j\arctan\bigl(k\tan\pi x_j\bigr)=\pi,
 \qquad \tau=-iy,\quad Q=e^y>1.
\end{equation}
A half label contributes $\pi/2$. The left side increases strictly
from zero or $\pi/2$ to $\pi\sum_jx_j$, so the solution exists,
is unique, and has a nonzero angular Hessian.

Let $V_m(\mathbf x)$ be the whole product action, after its original
hidden binary labels have been optimized. The exact elimination
identity is
\begin{equation}\label{gc:zero-legendre}
 V_m(\mathbf x)=\operatorname*{stat}_{\tau}
                  \left[\sum_j b(x_j,\tau)-i\tau\right].
\end{equation}
Both sides evaluate the same complete polynomial trace: one uses
the positive binary product, the other its first Fourier charge.
The constrained Laurent-coefficient calculation in
Appendix~\ref{app:strict-tree-geometric-characteristic} proves the
identity with its absolute constant. It continues holomorphically
through the possible half label.

The group marking requires a separate argument. In the fixed
product collar, the two based puncture loops satisfy
$a_{\mathrm{out}}=a_{\mathrm{in}}$ and
$b_{\mathrm{out}}=b_{\mathrm{in}}$, with disk-boundary word $ba$.
The complete old white-face maps identify the added primitive
longitude as $W_1^{-1}W_3$, up to its recorded orientation. Its
determinant-one normalization has trace $-(Q+Q^{-1})$, with
the same $Q>1$ as in \eqref{gc:zero-angular}. Pulling the old
holonomy back through the literal meridian-filling quotient keeps
both puncture generators and their paths. The resulting ordered
page conjugacies are the native all-circle conjugacies. Their
puncture pairs are irreducible, so the page intertwiners are unique
in the projective group. This identifies representations, including
the original odd seams, rather than matching only their traces.

At an internal half edge, $X_e=-1$ may be a double root of the
edge quadratic. The normalized character matrices remain regular.
One therefore retains the joint stationary system. The new angular
Hessian block is invertible; its Schur complement is exactly the
old strict reduced Hessian by \eqref{gc:zero-legendre}. The full
system is nondegenerate without inverting an individual edge
quadratic. Its rational longitude identity consequently extends
through this fold.

The selected added longitude remains semisimple near $Q>1$.
Its commuting meridian of eigenvalue ratio one is therefore
projectively identity, which removes the potential unipotent
ambiguity. The old marked deformation and sufficiently long
fillings identify the resulting main-meridian germ with that of
the original complete filled knot. Applying the all-circle
identity, and the same filled-knot Seifert sign, proves
\eqref{gc:strict-longitude}. No quantum estimate for the newly
drilled diagram is required; all quantum estimates remain those
of the original strict presentation.

\section{The resulting factor of a genuine quantum recurrence}
\label{gc:operator}

Now impose the all-order source hypotheses as well. They hold for
all the fixed strict Conway presentations in
Theorem~\ref{gc:strict-longitude-theorem} under independently
sufficiently long signed powers. For circle trees they also hold
under the summed condition
$\sum_vd_v/(d_v^2+n_v^2)\le2^{-25}$, or the local condition
$\max_vd_v/(d_v^2+n_v^2)\le2^{-120}$.
All original qualifications remain, and no data grow with the level.

Let $F_N(M)$ be the actual zero-framed reduced ADO state. A genuine
all-level annihilator means a fixed nonzero operator
\begin{equation}\label{gc:operator-definition}
 P(q,M,T)=\sum_{j=0}^d a_j(q,M)T^j,\qquad
 TM=qMT,\qquad
 P(q_N,M,T)F_N(M)=0
\end{equation}
at every sufficiently large level, where
$q_N=e^{\pi i/N}$ and $T$ shifts $M$ to $q_NM$.
The coefficients $a_j$ belong to $\mathbb C(q,M)$.
Clear its fixed rational denominators and remove the common
$q-1$ order. Setting $q=1$ and replacing the shift by a
commuting variable $L$ gives its first nonzero classical symbol
$P_0(M,L)$.

The all-order theorem in Appendix~\ref{app:filtered-characteristic}
uses the actual source on the translated microscopic charts
$s,s+1,\ldots,s+d$, away from ordinary resonances.
Its filtered logarithm selects the critical-value germ
$L_{\mathrm{eff}}=-e^{\mathcal S'(\log M)}$.
The exact edge insertion gives the right side of
\eqref{gc:strict-longitude}, with the original physical edge set.
The marked first-charge branch and the zero-circle elimination
fix the same filled representation. Theorems~\ref{gc:longitude-theorem}
and~\ref{gc:strict-longitude-theorem} therefore identify
$L_{\mathrm{eff}}$ with $L_{\mathrm{geom}}$ in their respective
source classes.

The reduced/full conversion retains its meridian factor: before
formal-domain cancellation the classical equation is
\begin{equation}\label{gc:meridian-factor}
 (M-M^{-1})^d P_0(M,L_{\mathrm{geom}}(M))=0.
\end{equation}
It follows that $P_0$ vanishes on the geometric germ.
Let $A_{\mathrm{geom},\mathbb C}$ be the irreducible Laurent
polynomial defining the smallest complex algebraic curve containing
this geometric eigenvalue germ. Then
\begin{equation}\label{gc:factor}
                     A_{\mathrm{geom},\mathbb C}\mid P_0.
\end{equation}
If $a_j\in\mathbb Q(q,M)$, one instead contracts the prime of
that germ to
$\mathbb Q[M^{\pm1},L^{\pm1}]$; the chosen complex component
need not itself be defined over $\mathbb Q$.

This is a selected geometric factor of every specified annihilator
under the source hypotheses. It does not identify the whole
classical symbol, the full A-polynomial, or the entire annihilator
ideal. Nor does it give a fixed-$M\ne1$ quantum asymptotic by
substituting $s$ of order $N$. The full geometric germ is obtained
from all orders on fixed microscopic charts and the separate
marked representation argument above.

\section{Transfer under an ordered reduction of the presentation}
\label{gc:intrinsic-transfer}

The reduced ADO function and the preferred eigenvalue germ belong
to the marked, zero-framed knot. Suppose the existing ordered
net-twist reduction gives an ambient isotopy from the input to a
covered strict presentation. The isotopy carries the knot
orientation, positive meridian, preferred zero-linking longitude
and cut mark. After correcting the transported ribbon framing to
zero, the full function $F_N(M)$ agrees at every level, as does
its rational continuation. Any genuine annihilator of that
intrinsic family therefore has the factor in \eqref{gc:factor}.
Presentation-specific recurrence formulas need not agree.

The same conclusion applies to an arbitrary-closure reduction
when it produces one covered strict knot factor and only unknot
factors. Under the reduction theorem's long surviving-factor
hypotheses, a hyperbolic knot input has precisely that form.
Removing the exhibited unknot summands is an ambient isotopy,
and its exterior map carries the preferred peripheral pair.
No assertion about quantum-growth multiplicativity for connected
sums is used here.

For a fixed original expression there are finitely many reduction
and parity branches. Taking the maximum of their strict source
and geometric thresholds gives a uniform sufficient threshold
for the surviving signed net powers. Large powers that cancel
to a short net power do not satisfy this condition. The integer
$a(D)$ and the edge factors are computed on the covered strict
diagram; only their preferred-longitude product and the selected
geometric factor are transferred intrinsically. The full proof
and precise reduction references are in
Section~\ref{gj:transfer-to-presentations-with-net-twists-and-arbitrary-closures}.
